\section{Feature parameterization and weight-space robustness}
\label{sec:features}

The benchmark payoffs in Section~\ref{sec:benchmarks} are generated by a transparent feature-based parameterization. This appendix makes every numerical input inspectable and provides a template for later calibration, but it does not identify the weights empirically or make the synthetic payoffs non-arbitrary.

\begin{definition}[Feature-based valuation]
Let $\Prof=A\times B$ be a finite profile set. A feature representation is a map
\[
\phi:\Prof\to\R^k.
\]
A context-specific weight matrix is a matrix
\[
W_C\in\R^{2\times k}.
\]
The induced contextual valuation is
\[
V_C(\sigma)=W_C\phi(\sigma).
\]
The first row of $W_C$ determines the actor-side payoff, and the second row determines the affected-side payoff.
\end{definition}

The same feature vector may be evaluated differently in different contexts. For example, public pressure may increase short-term compliance in a business organization, damage pedagogical legitimacy in a classroom, and improve focus in a sport team when embedded in a shared competitive culture.

\begin{proposition}[Feature weights induce contextual revaluation]
Let $\phi:\Prof\to\R^k$ be fixed. If $W_C\ne W_D$, then the same profile $\sigma$ may satisfy
\[
Q(W_C\phi(\sigma))\ne Q(W_D\phi(\sigma)).
\]
\end{proposition}

\begin{proof}
It suffices to give an example. Let $k=1$, $\phi(\sigma)=1$, $W_C=(1,-1)^T$ and $W_D=(1,1)^T$. Then $W_C\phi(\sigma)=(1,-1)$ is malicious, while $W_D\phi(\sigma)=(1,1)$ is intelligent. Hence changing the context weights can change the class while the profile feature vector is fixed.
\end{proof}

For the benchmark examples we use two abstract strategies for the authority-side player:
\[
G=\text{constructive guidance},\qquad P=\text{public pressure},
\]
and two abstract strategies for the affected-side player:
\[
U=\text{engaged uptake},\qquad R=\text{resistance or withdrawal}.
\]
The symbol $U$ is used for uptake in order to avoid confusion with the destructive class $\Dclass$.

The feature vector is
\[
\phi(\sigma)=(g,p,u,r,pu,gr,pr)^T\in\R^7,
\]
where $g,p,u,r$ are indicators for guidance, pressure, uptake, and resistance, while $pu,gr,pr$ are interaction indicators for $(P,U)$, $(G,R)$, and $(P,R)$ respectively. Table~\ref{tab:feature-meaning} explains the features. Table~\ref{tab:features} gives the four feature vectors.

\begin{table}[h]
\centering
\small
\begin{tabular}{c p{0.74\textwidth}}
\toprule
Feature & Interpretation\\
\midrule
$g$ & Baseline effect of constructive guidance by the authority-side player.\\
$p$ & Baseline effect of public pressure by the authority-side player.\\
$u$ & Baseline effect of engaged uptake by the affected-side player.\\
$r$ & Baseline effect of resistance or withdrawal by the affected-side player.\\
$pu$ & Contextual interaction effect when public pressure is met with engaged uptake.\\
$gr$ & Contextual interaction effect when guidance is met with resistance.\\
$pr$ & Contextual interaction effect when pressure is met with resistance.\\
\bottomrule
\end{tabular}
\caption{Interpretation of the benchmark features. The weights assigned to these features encode context-specific norms, goals, and role expectations.}
\label{tab:feature-meaning}
\end{table}

\begin{table}[h]
\centering
\small
\begin{tabular}{cccccccc}
\toprule
Profile & $g$ & $p$ & $u$ & $r$ & $pu$ & $gr$ & $pr$\\
\midrule
$(G,U)$ & 1 & 0 & 1 & 0 & 0 & 0 & 0\\
$(G,R)$ & 1 & 0 & 0 & 1 & 0 & 1 & 0\\
$(P,U)$ & 0 & 1 & 1 & 0 & 1 & 0 & 0\\
$(P,R)$ & 0 & 1 & 0 & 1 & 0 & 0 & 1\\
\bottomrule
\end{tabular}
\caption{Feature representation for the four benchmark profiles.}
\label{tab:features}
\end{table}

\begin{table}[h]
\centering
\small
\begin{tabular}{llrrrrrrr}
\toprule
Context & Payoff row & $g$ & $p$ & $u$ & $r$ & $pu$ & $gr$ & $pr$\\
\midrule
$C_B$ business & actor & 1 & 0 & 1 & -2 & 1 & 0 & 0\\
 & affected & 1 & -1 & 2 & -2 & -3 & 0 & 0\\
\midrule
$C_A$ classroom & actor & 1 & -2 & 1 & -3 & 0 & 0 & 3\\
 & affected & 1 & -4 & 2 & 0 & 0 & 0 & 1\\
\midrule
$C_S$ sport team & actor & 2 & 1 & 1 & -3 & 0 & 0 & 0\\
 & affected & 1 & 0 & 2 & -2 & 0 & 0 & -1\\
\bottomrule
\end{tabular}
\caption{Context-specific weight matrices. Applying $V_C(\sigma)=W_C\phi(\sigma)$ generates the benchmark payoff matrices in Table~\ref{tab:payoff-transitions}.}
\label{tab:weights}
\end{table}

The weights in Table~\ref{tab:weights} are illustrative. Their signs encode the intended institutional readings: in the business benchmark, pressure can increase short-term actor-side control but reduce the affected side's autonomy and trust; in the classroom benchmark it is penalized on both sides; and in the sport benchmark it is positive only when combined with uptake. These statements define the examples and are not empirical claims about the three domains.

The representation is deliberately saturated relative to the four observed profiles: seven features generate only four payoff vectors per context, so the weight matrix is not identified from those vectors. Multiple weight matrices produce the same game. Accordingly, the feature layer is an auditable parameterization rather than an explanatory regression model. The formal robustness claims in this appendix concern sign-defined sets of payoffs or weights; empirical identification would require additional profiles, restrictions, or external data.


\begin{proposition}[Class regions in weight space]
\label{prop:weight-space-regions}
Fix a profile $\sigma$ and its feature vector $\phi(\sigma)\in\R^k$. For each open class $c\in\{\Iclass,\Hclass,\Mclass\}$, the set of weights $W\in\R^{2\times k}$ such that $Q(W\phi(\sigma))=c$ is an intersection of two strict linear half-spaces. The strict destructive region $\mathcal D^-$ is also an intersection of two strict linear half-spaces. Thus, away from $\mathcal D^0$ and the boundary, each class condition is stable under sufficiently small perturbations of $W$.
\end{proposition}

\begin{proof}
Write the two rows of $W$ as $w_1,w_2\in\R^k$. Then
\[
W\phi(\sigma)=\bigl(w_1\cdot\phi(\sigma),\, w_2\cdot\phi(\sigma)\bigr).
\]
For fixed $\phi(\sigma)$, the maps $W\mapsto w_1\cdot\phi(\sigma)$ and $W\mapsto w_2\cdot\phi(\sigma)$ are linear functions of the entries of $W$. The class conditions are therefore linear sign conditions on these two functions. Explicitly,
\[
Q(W\phi(\sigma))=\Iclass \Longleftrightarrow w_1\cdot\phi(\sigma)>0\text{ and }w_2\cdot\phi(\sigma)>0,
\]
\[
Q(W\phi(\sigma))=\Hclass \Longleftrightarrow w_1\cdot\phi(\sigma)<0\text{ and }w_2\cdot\phi(\sigma)>0,
\]
\[
Q(W\phi(\sigma))=\Mclass \Longleftrightarrow w_1\cdot\phi(\sigma)>0\text{ and }w_2\cdot\phi(\sigma)<0,
\]
and
\[
W\phi(\sigma)\in\mathcal D^- \Longleftrightarrow w_1\cdot\phi(\sigma)<0\text{ and }w_2\cdot\phi(\sigma)<0.
\]
Each condition is the intersection of two strict linear half-spaces in $\R^{2\times k}$. Strict linear half-spaces are open sets, and finite intersections of open sets are open. Hence if a weight matrix $W$ places $W\phi(\sigma)$ in one of these open classes, there exists a radius $\rho>0$ such that every weight matrix $W'$ with $\|W'-W\|<\rho$ induces the same class for $\sigma$. The exceptional cases are precisely those in which one of the two linear forms is zero, corresponding to $\mathcal D^0$ or to the boundary of the benefit--loss partition.
\end{proof}

This proposition addresses sensitivity, not empirical validity. The benchmark tables are one point in weight space, and nearby weights preserving the induced signs leave the class conclusions unchanged. This does not establish that the selected region describes any real institution.

\subsection{Robustness of benchmark transitions in weight space}

The next result is the main formal response to the concern that the benchmark payoffs are hand-crafted. It shows that the reported class transitions are not isolated numerical coincidences: they persist on open polyhedral regions of the weight space. Thus the exact entries in Table~\ref{tab:weights} are illustrative representatives of broader sign-defined parameter regions.

\begin{proposition}[Open weight conditions for the benchmark transitions]
\label{prop:benchmark-weight-conditions}
Let $\phi_{PU}=\phi(P,U)$ and $\phi_{GR}=\phi(G,R)$. Write the rows of the business, classroom, and sport-team weight matrices as
\[
W_B=(w_B^A,w_B^B),\qquad W_A=(w_A^A,w_A^B),\qquad W_S=(w_S^A,w_S^B),
\]
where the superscripts $A$ and $B$ denote actor-side and affected-side rows. Then the transition
\[
(P,U):\Mclass\longrightarrow\Dclass\longrightarrow\Iclass
\]
holds, with the classroom point in $\mathcal D^-$, if and only if
\[
w_B^A\cdot\phi_{PU}>0,\quad w_B^B\cdot\phi_{PU}<0,
\]
\[
w_A^A\cdot\phi_{PU}<0,\quad w_A^B\cdot\phi_{PU}<0,
\]
and
\[
w_S^A\cdot\phi_{PU}>0,\quad w_S^B\cdot\phi_{PU}>0.
\]
Similarly,
\[
(G,R):\Dclass\longrightarrow\Hclass\longrightarrow\Dclass
\]
with the destructive points in $\mathcal D^-$ holds if and only if
\[
w_B^A\cdot\phi_{GR}<0,\quad w_B^B\cdot\phi_{GR}<0,
\]
\[
w_A^A\cdot\phi_{GR}<0,\quad w_A^B\cdot\phi_{GR}>0,
\]
and
\[
w_S^A\cdot\phi_{GR}<0,\quad w_S^B\cdot\phi_{GR}<0.
\]
Each set of inequalities defines an open polyhedral region in the product of the three weight spaces.
\end{proposition}

\begin{proof}
For any context $C$, the class of a profile is determined by the signs of the two coordinates of $W_C\phi(\sigma)$, namely
\[
W_C\phi(\sigma)=\bigl(w_C^A\cdot\phi(\sigma),\, w_C^B\cdot\phi(\sigma)\bigr).
\]
For $(P,U)$ to be malicious in the business context, these two signs must be $(+,-)$, giving the first pair of inequalities. For $(P,U)$ to be strictly destructive in the classroom context, the signs must be $(-,-)$, giving the second pair. For $(P,U)$ to be intelligent in the sport-team context, the signs must be $(+,+)$, giving the third pair. This proves the first equivalence.

The proof for $(G,R)$ is identical. The transition $\Dclass\to\Hclass\to\Dclass$ with strict destructive endpoints requires signs $(-,-)$ in business, $(-,+)$ in classroom, and $(-,-)$ in sport. Translating these sign requirements into dot-product inequalities gives the displayed conditions.

Finally, each displayed condition is a strict linear inequality in the entries of the corresponding weight row. A finite conjunction of strict linear inequalities is an open polyhedral region. Therefore the benchmark transitions are not isolated numerical coincidences; they persist throughout open regions of the relevant weight spaces.
\end{proof}

The inequalities also identify which signs are substantively essential. For example, the transition $(P,U):\Mclass\to\Dclass\to\Iclass$ requires public pressure with uptake to be actor-beneficial and affected-side harmful in the business context, harmful to both sides in the classroom context, and beneficial to both sides in the sport-team context. If any of these sign requirements is reversed, the transition changes. The feature weights should therefore be read as a compact representation of institutional norms and role expectations rather than as empirical estimates.

\begin{table}[h]
\centering
\small
\begin{tabular}{ccccc}
\toprule
Profile & Business $C_B$ & Classroom $C_A$ & Sport team $C_S$ & Class transition\\
\midrule
$(P,U)$ & $(+,-)$ & $(-,-)$ & $(+,+)$ & $\Mclass\to\Dclass\to\Iclass$\\
$(G,R)$ & $(-,-)$ & $(-,+)$ & $(-,-)$ & $\Dclass\to\Hclass\to\Dclass$\\
\bottomrule
\end{tabular}
\caption{Essential sign patterns behind the two recontextualized benchmark profiles. The exact numerical payoffs in Table~\ref{tab:payoff-transitions} are representatives of these sign-defined regions.}
\label{tab:essential-signs}
\end{table}

\begin{table}[h]
\centering
\small
\setlength{\tabcolsep}{4pt}
\begin{tabular}{>{\raggedright\arraybackslash}p{0.19\textwidth}
                >{\centering\arraybackslash}p{0.12\textwidth}
                >{\raggedright\arraybackslash}p{0.60\textwidth}}
\toprule
Context & Sign of $(P,U)$ & Institutional reading of the essential signs\\
\midrule
Business $C_B$ & $(+,-)$ & Public pressure met with uptake raises the actor's short-term control but erodes the affected side's autonomy and trust.\\
Classroom $C_A$ & $(-,-)$ & Public pressure damages pedagogical legitimacy and student confidence, harming both sides.\\
Sport team $C_S$ & $(+,+)$ & Public pressure accepted as tactical discipline benefits both the actor and the team.\\
\bottomrule
\end{tabular}
\caption{Institutional reading of the essential sign conditions for the recontextualized profile $(P,U)$. It is the signs, not the exact magnitudes, that carry the social interpretation; this is exactly what Proposition~\ref{prop:benchmark-weight-conditions} formalizes.}
\label{tab:institutional-signs}
\end{table}

\FloatBarrier
